\documentclass[../main.tex]{subfiles}
\begin{document}
%%%%%%%%%%%%%%%%%%%%%%%%%%%%%%%%%%%%%%%%%%%%%%%%%%%%%%%%%%%%%%%%%%%%%%%%%%%%%%%%%%%%%%%%%%%%%%%%%%%
\chapter{Using Templates and Diagram Styles}

When you open a new OmniGraffle document, it calls upon a \textit{template} that defines the
canvas, the units of measurement, and more. For example, the default template, Blank, has a
plain white canvas with the units set to pixels. The Blank template is simple, and it's a good
starting point for users new to OmniGraffle.

Conversely, a \textit{diagram style} is more of a snapshot of the styles used in an OmniGraffle
template that you can apply to your projects. Diagram styles allow you, for example, to
quickly prototype something without having to worry about applying styles as you go, and to
then sweep in at the end, apply the diagram style, and be done with your work.

\begin{ognote}
  The really cool part, though, is that a Diagram Style is nothing more than an OmniGraffle
  template file (\textit{.gtemplate}) that contains styled objects. When you apply a Diagram
  Style to a project, OmniGraffle references the objects in the template and applies matching
  styles to the current document.
\end{ognote}

This chapter introduces you to \textbf{Templates} and \textbf{Diagram Styles}, showing you how
to quickly apply styles from an existing template to help speed up your workflow.

\section{Using Templates}

In OmniGraffle, a \textit{template} is used to define the base properties of a project. When you
create a new document---by either choosing \textbf{File} $\blacktriangleright$ \textbf{New} or
pressing \textbf{Command-N}---OmniGraffle loads the default template into the window,
supplying you with a single canvas and its base layer.

As noted earlier, the default template for OmniGraffle 7 is the \textbf{Blank} template. It's
pretty sparse, which is great when you're just getting started. To choose a different
template, rather than choosing \textbf{File} $\blacktriangleright$ \textbf{New} to create a
new document, choose \textbf{File} $\blacktriangleright$ \textbf{Resource Browser} instead.
The Resource Browser is kind of like an OmniGraffle-specific Finder, that gives you access to
your OmniGraffle files, templates, and stencils.

\ogfigure{0.800}{res/og7_templates01.png}

After opening the Resource Browser, you can either click the \textbf{Templates} heading in the
sidebar to view all of the available templates, or select one of the folders for something
more specific to the project you plan to work on.

To learn more about a particular template, first click to select a template from one of the
folders, and then click \ogicon{12.5pt}{res/og7_button_rb_info01.png} in the Resource
Browser's upper-right corner. This button opens an Info pane along the right side of the
Resource Browser that displays additional information about the selected template.

\ogfigure{0.800}{res/og7_templates01a.png}

The Info pane displays a preview of the template, along with additional information about the
template's settings, including whether auto-layout is on or off, the dimensions of the canvas,
the units of measurement, and whether snap to grid is on or off.

\begin{ognote}
  After opening the Resource Browser's Info pane, that pane remains available until you close it
  again by clicking \ogicon{12.5pt}{res/og7_button_rb_info02.png}.
\end{ognote}

To open a new OmniGraffle project using the selected template, click \textbf{New Document}. Or,
if you'd like to change the default template to something else, click \textbf{Set as Default}.

\ogfigure{0.769}{res/og7_templates02.png}

When you change the default template, the area beneath the template preview lets you decide what
happens when you choose \textbf{File} $\blacktriangleright$ \textbf{New} (or press
\textbf{Command-N}) when creating a new document:

\ogfigure{0.769}{res/og7_templates03.png}

You can choose from:

\begin{itemize}
  \item \textbf{Open the Resource Browser}, which opens the Resource Browser window so you can
    choose a file, another template, or stencil to work on; or,
  \item \textbf{Use the Default Template}, which applies the default template to your
    OmniGraffle project when you click \textbf{New Document} in the Resource Browser's window.
\end{itemize}

\begin{ogtip}
  You can also set the default template in the Resource Browser by right-clicking on a template
  and choosing \textbf{Set as Default} from the contextual menu, or in OmniGraffle's General
  preferences by clicking the \textbf{Choose} button in the \textbf{New Documents} section.
\end{ogtip}

\subsection{Editing an Existing Template}

If you ever need to edit a template---say, to change turn on snap to grid, add ruler guides to
the canvas, or to add a background color---OmniGraffle makes it easy:

\begin{enumerate}
  \item Choose \textbf{File} $\blacktriangleright$ \textbf{Resource Browser}.
  \item \textbf{Control-/Right-click} the template that you want to edit, and then choose
    \textbf{Edit} if you want to edit and save the template file, or choose \textbf{Edit a Copy}
    from the menu:

    \ogfigure{0.800}{res/og7_templates04.png}

    \begin{itemize}
      \item If you choose \textbf{Edit}, the template file opens for you to work on. Once the
        template is open in OmniGraffle, you can work on it just like you would with any other
        project. When you have finished making your changes, saving the file overwrites the
        original file in the Templates folder.
      \item If you choose \textbf{Edit a Copy}, the contents of the template file are placed in a
        new, Untitled template file. When you've finished making your changes, be sure to give
        this template file a different name than its predecessor so you don't overwrite the
        original.
    \end{itemize}
\end{enumerate}

\subsection{Creating a New Template}

To create a new template, choose \textbf{File} $\blacktriangleright$ \textbf{New Resource}
$\blacktriangleright$ \textbf{New Template}, and then use the inspectors to set the properties
for the template. For example, you can:

\begin{itemize}
  \item Place Ruler Guides on the canvas so you don't have to recreate these for \textit{every}
    project that uses the template
  \item Turn on \textbf{Auto layout} in the Diagram Layout and Style inspector
  \item Choose a default line type in the Line inspector
  \item Specify an Infinite Canvas in the Canvas Size inspector
  \item Set a background color or an image to use on the canvas in the Canvas Fill or Background
    Image inspectors
  \item Enter information about the project file in the Document Data inspector
  \item and more!
\end{itemize}

When you save the template, OmniGraffle places the file inside the app container's
\textbf{Templates} folder, or you can select one of the subfolders (or create a new one) in
which to store your templates.

\ogfigure{0.800}{res/og7_templates05.png}

To view your template, choose \textbf{File} $\blacktriangleright$ \textbf{Resource Browser}
(\textbf{Shift-Command-N}), and choose the folder in which you saved your template. If you
didn't select a subfolder within the Templates folder, select the \textbf{Miscellaneous}
folder in the Resource Browser's sidebar, and then select your template.

\ogfigure{0.800}{res/og7_templates06.png}

If you want to set your newly-created template as the default, click \textbf{Set as Default},
and then choose one of the two options for what happens when you choose \textbf{File}
$\blacktriangleright$ \textbf{New}, as shown earlier.

\subsection{Creating a Shortcut to the Templates Folder}

To make the process of exporting an OmniGraffle file as a Template or Diagram Style a little
easier, you'll need to create a shortcut to the Templates folder in the Finder's sidebar.

The Templates folder lives deep within OmniGraffle's container, but there's a quick way to get
to the Stencils folder from within OmniGraffle. Follow along\ldots{}

\begin{enumerate}
  \item Choose \textbf{File} $\blacktriangleright$ \textbf{Resource Browser}, or use the
    \textbf{Shift-Command-N} keyboard shortcut to open the Resource Browser.
  \item In the Resource Browser, select any template, then click the Gear menu and choose
    \textbf{Show in Finder}:

    \ogfigure{0.800}{res/og7_templates_shortcut01.png}

    This opens a Finder window with that particular template selected:

    \ogfigure{0.800}{res/og7_templates_shortcut02.png}
  \item With the Finder window selected, choose \textbf{View} $\blacktriangleright$ \textbf{as
    Columns} (\textbf{Command-3}).
  \item Scroll the Finder view backward until you see the \textbf{Templates} folder; click to
    select the Templates folder:

    \ogfigure{0.800}{res/og7_templates_shortcut03.png}
  \item With the Templates folder selected, choose \textbf{File} $\blacktriangleright$
    \textbf{Add to Sidebar} (\textbf{Control-Command-T}). This adds a shortcut to the Templates
    folder in the \textbf{Favorites} section of the Finder's sidebar.

    \ogfigure{0.740}{res/og7_templates_shortcut04.png}
\end{enumerate}

Now when you need to export an OmniGraffle file to use as a Template or Diagram Style, you can
choose the \textbf{Templates} folder in the Favorites list on the left side of the save sheet.

\ogfigure{0.800}{res/og7_templates_shortcut05.png}

The other benefit of knowing where the Templates folder is on your Mac is so you can quickly add
subfolders to organize your Templates and Diagram Styles, as well as rename or delete
templates, too.

\begin{ogsidebar}
  If you need to access the Templates folder from a script you are writing, the Unix path to use
  is:

  \begin{itemize}
    \item \textbf{Purchased from the Omni Group:}

      \noindent\footnotesize\path{~/Library/Containers/com.omnigroup.OmniGraffle7/Data/Library/Application\ Support/The\ Omni\ Group/OmniGraffle/Templates}
    \item \textbf{Purchased from the Mac App Store:}

      \noindent\footnotesize\path{~/Library/Containers/com.omnigroup.OmniGraffle7.MacAppStore/Data/Library/Application\ Support/The\ Omni\ Group/OmniGraffle/Templates}
  \end{itemize}

  Triple-click the desired path, press \textbf{Command-C} to copy the path to the Clipboard, and
  then paste (\textbf{Command-V}) it into your code editor.
\end{ogsidebar}

\section{Using Diagram Styles}

Diagram Styles, which you can apply to a diagram using the Diagram Layout and Style inspector,
are nothing more than an OmniGraffle Template file (\textit{.gtemplate}) with styled objects
on the canvas. The difference between a \textit{template} and a \textit{diagram style},
however, is that templates are used as the basis for a new OmniGraffle project, while a
diagram style is \textit{applied} to an OmniGraffle project that you are already working on.

Still confused? Don't be. Let's explore this further\ldots{}

\begin{enumerate}
  \item Choose \textbf{File} $\blacktriangleright$ \textbf{Resource Browser}
    (\textbf{Shift-Command-N})
  \item In the sidebar under \textbf{Templates}, select \textbf{Diagrams} and then select the
    \textbf{Hierarchical Lines} template:
\end{enumerate}

\begin{ogwarning}
  Do not click \textbf{New Document} quite yet. Keep following to learn more about Diagram
  Styles\ldots{}
\end{ogwarning}

\ogfigure{0.800}{res/og7_diagramstyles01.png}

As you can see by looking at the Hierarchical Lines template's preview in the Resource Browser,
this particular template has a black background and a tree of colored lines with labels to
denote a diagram's hierarchy. If you click \textbf{New Document}, OmniGraffle opens a new
Untitled document window using the contents of the Hierarchical Lines template, including the
diagram that is part of the template.

\ogfigure{0.800}{res/og7_diagramstyles02.png}

Now, you might be thinking, \textit{``Why would I want to open a project with stuff already on
the canvas?''}. You probably wouldn't want to, and there's nothing that's stopping you from
selecting everything that's on the canvas, pressing \textbf{Delete} and starting with a clean
slate. But\ldots{}

\ldots{}there may be times when you actually \textit{want} the stuff that's already on the
canvas. But more importantly, when you create a new document that uses one of the templates in
the Diagrams folder, what you are essentially getting is a pre-styled diagram. However, now
you have to change all of the labels, and that kind of work can be tedious.

Instead, what you really want is to \textit{apply} the styles from one of the templates in the
Diagrams folder to a diagram that you are already working on. Here's how you do that:

\begin{enumerate}
  \item Choose \textbf{File} $\blacktriangleright$ \textbf{New} from the menu bar.
  \item Use the \textbf{Outline} tab in the Sidebar to create the structure for your diagram.

    \ogfigure{0.800}{res/og7_diagramstyles03.png}

    The diagram in its current state doesn't really look like much; certainly not something that
    would impress your boss or friends. At this point, you have two options:
  \item Style objects individually (see Diagramming Basics), which could take a little while;
    or,
  \item Apply a \textbf{Diagram Style} to the diagram.

    Let's go with Option 2.
  \item Switch to the Canvas Inspector pane and navigate to the Diagram Layout and Style
    inspector. Click \textbf{Choose\ldots{}} next to \textbf{Style} in the inspector. This opens
    a sheet displaying the available Templates in the left sidebar.

    \ogfigure{0.800}{res/og7_diagramstyles11.png}
  \item Select the \textbf{Diagrams} folder in the left sidebar, and then click to select the
    \textbf{Hierarchical Lines} template.

    \ogfigure{0.800}{res/og7_diagramstyles04.png}
  \item Click \textbf{Apply} to apply the styles from the Hierarchical Lines template to the
    objects in the diagram:

    \ogfigure{0.800}{res/og7_diagramstyles05.png}
\end{enumerate}

When you apply a diagram style, the items of the diagram are styled to match the template's
styles for shapes, connection lines, text objects, and automatic layout settings. If your
diagram has more levels than in the template you have applied, the deeper items use the styles
of the deepest level of the template's style.

\subsection{Customizing an Existing Template's Diagram Styles}

Now that you know how to apply a diagram style to your projects, the final step in this journey
is to customize an existing template's diagram styles to create one that's more \textit{your}
style.

\begin{enumerate}
  \item Start by opening the Resource Browser (\textbf{Shift-Command-N}), and choose
    \textbf{Diagrams} in the Templates section of the sidebar:

    \ogfigure{0.800}{res/og7_diagramstyles06.png}
  \item \textbf{Control-/Right-click} on the template that you would like to customize, and
    choose \textbf{Edit a Copy} from the contextual menu that appears:

    \ogfigure{0.800}{res/og7_diagramstyles07.png}

    This opens a copy of the template file so you can change things up and then save it as a
    different file name.

    \ogfigure{0.800}{res/og7_diagramstyles08.png}
  \item Change the styles. For example, with the Hierarchical Lines template, you could change
    the background color, the font and text styles used for the object labels, balance out the
    object styles so everything at the same hierarchical level has the same style properties,
    and more.
  \item When you have finished, save the template as a different filename. Templates are saved
    by default to the Templates folder; however, you can click the \textbf{Where} pop-up menu to
    select a subfolder within the Templates folder (such as the \textbf{Diagrams} folder), if
    you'd like.

    \ogfigure{0.800}{res/og7_diagramstyles09.png}
\end{enumerate}

You can double-check your work by going back to the Resource Browser and selecting the folder in
which you saved the template, or you can put it to use by applying the template's diagram
styles to an existing project.

\ogfigure{0.800}{res/og7_diagramstyles10.png}
%%%%%%%%%%%%%%%%%%%%%%%%%%%%%%%%%%%%%%%%%%%%%%%%%%%%%%%%%%%%%%%%%%%%%%%%%%%%%%%%%%%%%%%%%%%%%%%%%%%
\end{document}
